% ================================================================
% APPENDIX
% ================================================================

\section{Proofs}\label{app:proofs}

% TODO: Include proofs for:
%   - Proposition: decomposition of covariance into firm-level contributions

\subsection{Proof of Proposition \ref{prop: abatement}}

We prove the result in three steps: first, we show that non-targeted firms do not abate; second, we show that for targeted firms the first-order condition yields a smooth interior solution at $p_z = 0$; third, we compute the semi-elasticity of emission intensity with respect to $p_z$.

\paragraph{Step 1: Non-targeted firms.} Consider a firm with $\tau_i = 0$. Its cost minimization problem is
\begin{equation*}
    mc_i = \min_{\{x_{ij}\}, \ell_i, \kappa_i} \left\{ \sum_{j \in \N} p_j x_{ij} + w\ell_i + \frac{\kappa_i^2}{2} \quad \text{s.t.} \quad f_i(\{x_{ij}\}, \ell_i) \geq 1 \right\}
\end{equation*}
where the emission cost term $\tau_i p_z e_i$ vanishes. Since the objective is increasing in $\kappa_i$ for $\kappa_i > 0$ and $\kappa_i$ does not enter the production constraint (Assumption 1), the minimum is achieved at $\kappa_i^* = 0$. Thus, $e_i^* = \overline{e}_i$ and $d \log e_i / d p_z = 0$.

\paragraph{Step 2: Interior solution for targeted firms.} Consider a targeted firm ($\tau_i = 1$). Substituting the abatement function $e_i = (1-\kappa_i)^{\alpha_i}\overline{e}_i$ into the objective, the terms involving $\kappa_i$ are
\begin{equation*}
    p_z (1-\kappa_i)^{\alpha_i}\overline{e}_i + \frac{\kappa_i^2}{2}
\end{equation*}
Since $\kappa_i$ does not enter the constraint $f_i \geq 1$ (Assumption 1), the first-order condition for $\kappa_i$ is
\begin{equation}
    \label{eq: foc kappa}
    F(\kappa_i, p_z) := p_z \alpha_i (1-\kappa_i)^{\alpha_i - 1}\overline{e}_i - \kappa_i = 0
\end{equation}
At $p_z = 0$, this is satisfied by $\kappa_i^* = 0$: both the marginal benefit of abatement ($0 \cdot \alpha_i \overline{e}_i = 0$) and the marginal cost ($\kappa_i = 0$) are zero. This is an interior solution (tangency), not a corner.

To verify that this characterizes a minimum, note that the second derivative with respect to $\kappa_i$ is
\begin{equation*}
    p_z \alpha_i (\alpha_i - 1)(1-\kappa_i)^{\alpha_i - 2}\overline{e}_i + 1 > 0
\end{equation*}
which is strictly positive for all $(\kappa_i, p_z)$ with $\kappa_i \in [0,1)$.

\paragraph{Step 3: Semi-elasticity at $p_z = 0$.} We apply the implicit function theorem to (\ref{eq: foc kappa}). At $(\kappa_i, p_z) = (0, 0)$:
\begin{align*}
    \frac{\partial F}{\partial p_z}\bigg|_{(0,0)} &= \alpha_i (1 - 0)^{\alpha_i - 1} \overline{e}_i = \alpha_i \overline{e}_i \\
    \frac{\partial F}{\partial \kappa_i}\bigg|_{(0,0)} &= -0 \cdot \alpha_i(\alpha_i - 1)(1-0)^{\alpha_i - 2}\overline{e}_i - 1 = -1 \neq 0
\end{align*}
Since $\partial F / \partial \kappa_i \neq 0$, the implicit function theorem guarantees that $\kappa_i^*(p_z)$ is a smooth function of $p_z$ in a neighborhood of $p_z = 0$, with
\begin{equation}
    \label{eq: dkappa dpz}
    \frac{d\kappa_i^*}{dp_z}\bigg|_{p_z = 0} = -\frac{\partial F / \partial p_z}{\partial F / \partial \kappa_i}\bigg|_{(0,0)} = -\frac{\alpha_i \overline{e}_i}{-1} = \alpha_i \overline{e}_i
\end{equation}

Now, $\log e_i = \alpha_i \log(1 - \kappa_i) + \log \overline{e}_i$, so
\begin{equation*}
    \frac{d \log e_i}{d p_z} = -\frac{\alpha_i}{1 - \kappa_i^*} \cdot \frac{d\kappa_i^*}{dp_z}
\end{equation*}
Evaluating at $p_z = 0$ (where $\kappa_i^* = 0$) and substituting (\ref{eq: dkappa dpz}):
\begin{equation*}
    \frac{d \log e_i}{d p_z}\bigg|_{p_z = 0} = -\alpha_i \cdot \alpha_i \overline{e}_i = -\alpha_i^2 \overline{e}_i
\end{equation*}
Including the targeting dummy, the result for any firm $i$ is $d \log e_i / dp_z|_{p_z=0} = -\tau_i \alpha_i^2 \overline{e}_i$. \qed

\subsection{Proof of Proposition \ref{prop: price propagation}}

From $p_i = mc_i$ and the envelope theorem, the total derivative of the unit cost function with respect to $p_z$ is
\begin{equation*}
    \frac{dp_i}{dp_z} = \sum_{j \in \N} \frac{\partial mc_i}{\partial p_j}\frac{dp_j}{dp_z} + \frac{\partial mc_i}{\partial w}\frac{dw}{dp_z} + \frac{\partial mc_i}{\partial p_z}
\end{equation*}
By Shephard's lemma, $\partial mc_i / \partial p_j = x_{ij}^*$ and $\partial mc_i / \partial w = \ell_i^*$. For the direct effect of $p_z$, the envelope theorem gives $\partial mc_i / \partial p_z = \tau_i e_i^*$, since the indirect effects through the optimal choices $\kappa_i^*$, $x_{ij}^*$, and $\ell_i^*$ vanish at the optimum. Thus:
\begin{equation}
    \label{eq: total deriv p}
    \frac{dp_i}{dp_z} = \sum_{j \in \N} x_{ij}^* \frac{dp_j}{dp_z} + \ell_i^* \frac{dw}{dp_z} + \tau_i e_i^*
\end{equation}
Dividing both sides by $p_i$ and multiplying and dividing each term by the appropriate price to form cost shares:
\begin{equation*}
    \frac{d \log p_i}{dp_z} = \sum_{j \in \N} \frac{p_j x_{ij}^*}{p_i} \frac{d \log p_j}{dp_z} + \frac{w \ell_i^*}{p_i} \frac{d \log w}{dp_z} + \frac{\tau_i e_i^*}{p_i}
\end{equation*}
Evaluating at $p_z = 0$, where $e_i^* = \overline{e}_i$ and $p_i = 1$ (normalization), this becomes the forward equation:
\begin{equation}
    \label{eq: forward eq proof}
    \frac{d \log p_i}{dp_z} = \sum_{j \in \N} \Omega_{ij} \frac{d \log p_j}{dp_z} + \Omega_{iL} \frac{d \log w}{dp_z} + \tau_i \overline{e}_i
\end{equation}

In matrix form, letting $\mathbf{v} := d \log \mathbf{p} / dp_z$ denote the vector of price semi-elasticities for the $N$ goods (excluding the consumption and labor sectors):
\begin{equation*}
    \mathbf{v} = \Omega_{\N \times \N} \, \mathbf{v} + \Omega_{(:, L)} \frac{d \log w}{dp_z} + \tau \circ \overline{\mathcal{E}}
\end{equation*}
where $\Omega_{\N \times \N}$ is the $N \times N$ submatrix of intermediate input shares. Rearranging:
\begin{equation*}
    (I - \Omega_{\N \times \N}) \, \mathbf{v} = \tau \circ \overline{\mathcal{E}} + \Omega_{(:, L)} \frac{d \log w}{dp_z}
\end{equation*}
Since $(I - \Omega_{\N \times \N})$ is invertible with inverse $\Psi$ (the Leontief inverse), we obtain:
\begin{equation*}
    \mathbf{v} = \Psi (\tau \circ \overline{\mathcal{E}}) + \Psi_{(:,L)} \frac{d \log w}{dp_z}
\end{equation*}

Finally, with nominal output as numeraire and fixed labor supply ($d \log L = 0$), $d \log w = d \log \lambda_L$, yielding the result:
\begin{equation*}
    \frac{d \log p_i}{dp_z} = \sum_{j \in \N} \Psi_{ij} \tau_j \overline{e}_j + \Psi_{iL} \frac{d \log \lambda_L}{dp_z}
\end{equation*}
\qed

\subsection{Proof of Proposition \ref{prop: agg output doesn't move}}

We show that $d \log Y / dp_z|_{p_z = 0} = 0$ by establishing that the competitive equilibrium at $p_z = 0$ is efficient, and that the introduction of a carbon tax is a first-order-irrelevant perturbation of the allocation.

\paragraph{Step 1: The no-policy equilibrium is efficient.} At $p_z = 0$, the economy has no distortions: prices equal marginal cost ($p_i = mc_i$), profits are zero (CRS), tax revenue is zero, and abatement costs are zero ($\kappa_i^* = 0$ implies $g(\kappa_i^*) = 0$). By the first welfare theorem, the competitive equilibrium allocation $\{x_{ij}^*, \ell_i^*, y_i^*\}$ is Pareto efficient. With a representative household, this means the equilibrium allocation maximizes real output:
\begin{equation*}
    Y^* = \max_{\{x_{ij}, \ell_i, y_i\}} \mathcal{C}(y_1, \ldots, y_N) \quad \text{s.t.} \quad f_i(\{x_{ij}\}, \ell_i) \geq y_i + \sum_{j \in \N} x_{ji} \;\; \forall i, \quad \sum_{i \in \N} \ell_i \leq L
\end{equation*}

\paragraph{Step 2: The carbon tax does not change the feasible set.} Note that neither $p_z$, nor $\tau$, nor $\kappa_i$ appear in the resource constraints above. The production functions $f_i$ are unchanged by the carbon tax (Assumption 1: emission intensity is not a function of the input bundle). The carbon tax and abatement affect equilibrium prices and allocations, but the set of physically feasible allocations is independent of $p_z$.

\paragraph{Step 3: First-order irrelevance.} For $dp_z > 0$, the competitive equilibrium allocation changes: firms substitute across inputs and invest in abatement. Denote the new equilibrium allocation as a function of $p_z$. The resulting real output satisfies $Y(p_z) \leq Y^*$ for all $p_z$, since the distorted equilibrium is generically inefficient. Since $Y(0) = Y^*$ and $Y(p_z) \leq Y^*$, the function $Y(p_z)$ achieves its maximum at $p_z = 0$, and therefore:
\begin{equation*}
    \frac{d \log Y}{dp_z}\bigg|_{p_z = 0} = 0
\end{equation*}

Note that the convexity of the abatement cost ($g(\kappa_i) = \kappa_i^2/2$) plays a role here: since $g'(0) = 0$, the total abatement expenditure $\sum_i g(\kappa_i^*) x_i$ is second-order in $dp_z$ (because $\kappa_i^* \propto dp_z$ and $g(\kappa_i^*) \propto dp_z^2$). With a linear abatement cost, abatement expenditure would be first-order in $dp_z$, and one would need the additional assumption that abatement resources are rebated to the household for the result to hold. \qed

\subsection{Proof of Proposition \ref{prop: agg emissions}}

Start from $Z = \sum_{i \in \N} e_i x_i$. By the product rule:
\begin{equation*}
    d \log Z = \frac{dZ}{Z} = \sum_{i \in \N} \frac{e_i x_i}{Z}\left(\frac{de_i}{e_i} + \frac{dx_i}{x_i}\right) = \sum_{i \in \N} \frac{z_i}{Z}\left(d \log e_i + d \log x_i\right)
\end{equation*}

From the definition of the Domar weight $\lambda_i = p_i x_i / Q$ and the normalization that nominal output $Q$ is the numeraire ($d \log Q = 0$):
\begin{equation*}
    d \log \lambda_i = d \log p_i + d \log x_i - d \log Q = d \log p_i + d \log x_i
\end{equation*}
so that $d \log x_i = d \log \lambda_i - d \log p_i$. Substituting:
\begin{equation*}
    d \log Z = \sum_{i \in \N} \frac{z_i}{Z} d \log e_i + \sum_{i \in \N} \frac{z_i}{Z}\left[d \log \lambda_i - d \log p_i\right]
\end{equation*}
Dividing both sides by $dp_z$ and evaluating at $p_z = 0$ yields the result. The scale channel --- the change in aggregate emissions due to a change in aggregate output holding emission intensities and the allocation constant --- is absent because aggregate output does not move at first order (Proposition \ref{prop: agg output doesn't move}). \qed

\subsection{Proof of Theorem \ref{theorem}}

The proof proceeds in three steps. We first derive the equilibrium quantity responses in the CES economy, then expand the covariance form using structural identities, and finally show they coincide.

\paragraph{Step 1: Equilibrium quantity responses.}
In the CES economy, firm $i$'s cost of producing one unit of output decomposes as $p_i = c_i^{int} + \tau_i p_z e_i + \kappa_i^2/2$, where $c_i^{int} = (\sum_j \theta_{ij}p_j^{1-\sigma} + \theta_{i\ell}w^{1-\sigma})^{1/(1-\sigma)}$ is the CES input cost index. Cost minimization gives per-unit input demands with elasticity $\sigma$ with respect to relative input prices:
\begin{equation}
    \label{eq: per unit demand}
    \frac{d \log \hat{x}_{ij}}{dp_z} = -\sigma\left(\frac{d \log p_j}{dp_z} - \frac{d \log c_i^{int}}{dp_z}\right)
\end{equation}
where $\hat{x}_{ij}$ denotes per-unit demand (the amount of input $j$ needed to produce one unit of good $i$). By Shephard's lemma on $c_i^{int}$ and the forward equation (Proposition \ref{prop: price propagation}):
$$\frac{d \log c_i^{int}}{dp_z} = \sum_{m \in \U} \Omega_{im} v_m = v_i - \tau_i \overline{e}_i$$
where $v_m := d\log p_m / dp_z$. Total demands are $x_{ij} = \hat{x}_{ij} \cdot x_i$, so:
\begin{equation}
    \label{eq: total demand ces}
    \frac{d \log x_{ij}}{dp_z} = \frac{d \log x_i}{dp_z} - \sigma(v_j - v_i + \tau_i\overline{e}_i)
\end{equation}

For the household, $\tau_C = 0$ and $\overline{e}_C = 0$, and $d\log P/dp_z = 0$ (Proposition \ref{prop: agg output doesn't move} and the numeraire normalization), so final demand satisfies $d \log y_j / dp_z = -\sigma v_j$.

Aggregating through market clearing ($x_j = y_j + \sum_{i \in \N} x_{ij}$) and defining $\widetilde{B}_{ji} := x_{ij}/x_j$ (the fraction of $j$'s output sold to firm $i$ as intermediate input, with $\sum_{i \in \N}\widetilde{B}_{ji} < 1$ since $y_j > 0$):
\begin{equation}
    \label{eq: quantity system}
    (I - \widetilde{B})\mathbf{q} = -\sigma(I - \widetilde{B})\mathbf{v} - \sigma\widetilde{B}(\tau \circ \overline{\mathcal{E}})
\end{equation}
where $\mathbf{q} := d \log \mathbf{x}/dp_z$. Since $\widetilde{B}$ is sub-stochastic, $I - \widetilde{B}$ is invertible and:
\begin{equation}
    \label{eq: q solved}
    \mathbf{q} = -\sigma\mathbf{v} - \sigma(\widetilde{\Psi} - I)(\tau \circ \overline{\mathcal{E}})
\end{equation}
where $\widetilde{\Psi} := (I - \widetilde{B})^{-1}$ is the demand-side Leontief inverse, satisfying $\widetilde{\Psi} = D^{-1}\Psi^T D$ with $D = \text{diag}(\lambda_1, \ldots, \lambda_N)$.

\paragraph{Step 2: Expand the reallocation channel.}
From Proposition \ref{prop: agg emissions}, the reallocation channel is $R = \sum_{j \in \N}(z_j/Z)\, q_j$. Substituting (\ref{eq: q solved}) and using $z_j = \overline{e}_j x_j$, $\widetilde{\Psi}_{jk} = (\lambda_k/\lambda_j)\Psi_{kj}$, and $x_j = \lambda_j Q$:
\begin{equation}
    \label{eq: R expanded}
    R = -\frac{\sigma}{Z}\bigg[\underbrace{\sum_j \overline{e}_j x_j v_j}_{(A)} + \underbrace{\sum_j \overline{e}_j \sum_k x_k \Psi_{kj}\tau_k\overline{e}_k}_{(B)} - \underbrace{\sum_j \tau_j \overline{e}_j^2 x_j}_{(C)}\bigg]
\end{equation}

\paragraph{Step 3: Expand the covariance form and match.}
For any buyer $i \in \N \cup \{C\}$, we substitute the structural identities
\begin{align*}
    \sum_{m \in \U} \Omega_{im} v_m &= v_i - \tau_i\overline{e}_i & &\text{(forward equation)} \\
    \sum_{m \in \U} \Omega_{im} (\Psi_e)_m &= (\Psi_e)_i - \overline{e}_i & &\text{(Leontief identity: } \Psi_e = \overline{\mathcal{E}} + \Omega\Psi_e\text{)}
\end{align*}
into the covariance $\mathbb{C}_{\Omega_{(i,:)}}(v, \Psi_e) = \sum_m \Omega_{im}v_m(\Psi_e)_m - (v_i - \tau_i\overline{e}_i)((\Psi_e)_i - \overline{e}_i)$. Summing over all buyers weighted by $x_i/Z$:

\textit{Cross-moment term}: using market clearing, $\sum_{i \in \N \cup \{C\}} x_i \Omega_{im} = x_m$ for all $m \in \N$, so
$$\frac{1}{Z}\sum_{i \in \N \cup \{C\}} x_i \sum_{m} \Omega_{im}v_m(\Psi_e)_m = \frac{1}{Z}\sum_m x_m v_m (\Psi_e)_m$$

\textit{Mean-product term}: for the consumption sector, $\tau_C = \overline{e}_C = v_C = 0$, so only $i \in \N$ contributes. Expanding:
$$\frac{1}{Z}\sum_{i \in \N} x_i(v_i - \tau_i\overline{e}_i)((\Psi_e)_i - \overline{e}_i) = \frac{1}{Z}\bigg[\sum_i x_i v_i(\Psi_e)_i - \sum_i x_i v_i \overline{e}_i - \sum_i x_i\tau_i\overline{e}_i(\Psi_e)_i + \sum_i x_i\tau_i\overline{e}_i^2\bigg]$$

\textit{Combining}: the cross-moment and the first term of the mean-product cancel ($\sum_m x_m v_m(\Psi_e)_m = \sum_i x_i v_i(\Psi_e)_i$, same sum relabeled; the $C$ term vanishes since $v_C = 0$). We are left with:
$$\sum_{i \in \N \cup \{C\}} \frac{x_i}{Z}\mathbb{C}_{\Omega_{(i,:)}}(v, \Psi_e) = \frac{1}{Z}\sum_j \overline{e}_j x_j v_j + \frac{1}{Z}\sum_j \tau_j\overline{e}_j x_j\big[(\Psi_e)_j - \overline{e}_j\big]$$

\textit{Matching with (\ref{eq: R expanded})}: term $(A)$ matches the first sum. For $(B)$, swap the order of summation:
$$\sum_j \overline{e}_j \sum_k x_k\Psi_{kj}\tau_k\overline{e}_k = \sum_k \tau_k\overline{e}_k x_k \underbrace{\sum_j \Psi_{kj}\overline{e}_j}_{= (\Psi_e)_k} = \sum_k \tau_k\overline{e}_k x_k(\Psi_e)_k$$
Relabeling $k \to j$, this equals $\sum_j \tau_j\overline{e}_j x_j(\Psi_e)_j$, so $(B) - (C) = \sum_j \tau_j\overline{e}_j x_j[(\Psi_e)_j - \overline{e}_j]$, matching the second sum. \qed


\section{Data Appendix}\label{app:data}

% ------------------------------------------------------------------
% Sample coverage
% ------------------------------------------------------------------
\subsection{Sample Coverage}

% TODO: Include the sample coverage table showing for selected years:
%   - Number of firms in sample
%   - Number of EU ETS firms
%   - GDP (value added) coverage as % of aggregate
%   - Wage bill coverage as % of aggregate
%   - Emissions coverage as % of aggregate and as % of EUETS sectors

% ------------------------------------------------------------------
% Imputation details
% ------------------------------------------------------------------
\subsection{Imputation Details}

% TODO: Detailed description of the emissions imputation procedure:
%   - Data sources for sector-level emissions
%   - Concordance between energy balance sectors and NACE codes
%   - Validation exercises beyond the R^2 = 0.56 reported in the text
%   - Robustness to alternative imputation methods

% ------------------------------------------------------------------
% Covariance operator
% ------------------------------------------------------------------
\subsection{Covariance Operator}

% TODO: Formal definition of the input-weighted covariance operator:
%
% Definition: The input-weighted covariance operator between vectors
% A_{(:,m)} and B_{(:,n)} is
%
%   C_{Omega(i,:)}(A_{(:,m)}, B_{(:,n)}) :=
%     sum_j Omega_{(i,j)} A_{(j,m)} B_{(j,n)}
%     - (sum_j Omega_{(i,j)} A_{(j,m)}) (sum_j Omega_{(i,j)} B_{(j,n)})
%
% where Omega_{(i,:)} denotes the i-th row of Omega and A_{(:,m)} the
% m-th column of A. It is the covariance between A and B using input
% shares as probability weights.


\section{Additional Figures and Tables}\label{app:additional}

% TODO: Additional figures and tables referenced in the text.
